\subsection*{Terminal A.11/A.19 residuals as exact stress-gauge defects (22 October 2026)}
\label{subsec:A11-source-terminal-residual-lift}
Let \(\mathcal H\) be the full A.11 two-support hull and intervening gap.
On the source region \(U=M=J=0\), write
\[
 N_s=\frac{4h\eta S-(1-\eta^2)S_\eta}{X}
       +4A\eta\Pi-(1-\eta^2)\Pi_\eta,\quad A=\tfrac12+h.
\]
Fix the \emph{same} corrected source profile throughout and define
\[
 \mathcal F(\eta)=\frac{\lambda S(\infty,\eta)}
  {X_pe_b^2(1+\eta^2)^{-2}},\qquad
 \mathcal G(\eta)=\int_{\text{two A.11 bumps}}(E_c^2-E_b^2)\,dy.
\]
The source baseline pressure trace is chosen to cancel the baseline
terminal pressure. Thus, with
\(K(\eta)=X_pe_b^2/[\lambda(1+\eta^2)^2]\),
the genuine terminal errors are \(C_S=K\mathcal F\),
\(C_\Pi=\mathcal G/2\). Differentiating \emph{totally} in \(\eta\)
along the actual A.19 amplitude and A.11 coefficient branches gives
the exact identity
\begin{equation}\label{eq:a11-source-residual-lift-20261022}
\boxed{\begin{aligned}
 \Delta N_s={}&\frac{K}{X}
  \left[\left(4h\eta+
   \frac{4\eta(1-\eta^2)}{1+\eta^2}\right)\mathcal F
          -(1-\eta^2)\mathcal F_\eta\right]\\
 &+2A\eta\mathcal G-\frac{1-\eta^2}{2}\mathcal G_\eta.
\end{aligned}}
\end{equation}
This follows from \(K_\eta/K=-4\eta/(1+\eta^2)\).
Disjoint A.11 bumps with eta-independent baseline give
\(\mathcal G=\sum_{i=1}^2(2v_ic_i+w_ic_i^2)\) and
\(\mathcal G_\eta=2\sum_i(v_i+w_ic_i)c_{i,\eta}\),
with \(v_i=\int E_b^2\beta_i\,dy\), \(w_i=\int E_b^2\beta_i^2\,dy\).
Consequently, \emph{exact} \(C^1\) simultaneous solutions of the
A.19 and A.11 source equations imply
\(\mathcal F=\mathcal F_\eta=\mathcal G=\mathcal G_\eta=0\);
small residual values alone do not control their total derivatives.

\begin{proposition}[Fail-closed two-support-and-gap stress gate]
Suppose the authentic zero-gauge estimate on the \emph{same}
corrected profile is \(b_0\ge\sup_{\mathcal H}|N_s^{(0)}|/E\).
Assume \(E>0\) and that independently authenticated nonnegative
source bounds satisfy
\[
 \alpha\ge\sup_{\mathcal H}\frac{K}{XE},\qquad
 \beta\ge\sup_{\mathcal H}\frac1E,\qquad
 \epsilon_F\ge\|\mathcal F\|_\infty,\quad
 \epsilon_{F'}\ge\|\mathcal F_\eta\|_\infty,
\]
and analogously \(\epsilon_G,\epsilon_{G'}\).
For \(0<\lambda\le1/20,\;0<h\le\lambda/100\), one has
\[
 \sup_{\mathcal H}\frac{|N_s|}{E}\le b_0+
 \alpha[(2+4h)\epsilon_F+\epsilon_{F'}]
 +\beta[2A\epsilon_G+\epsilon_{G'}/2].
\]
If \(b_0\le\lambda/500\) and the two residual brackets total
strictly less than \(249\lambda/500\), then
\(\sup_{\mathcal H}|N_s|/E<\lambda/2\).
\end{proposition}
\begin{proof}
Use \(|\eta|\le1\), \(0\le1-\eta^2\le1\), and the exact identity
\(2(1+u^2)-4u(1-u^2)=2[(1-u)^2+2u^3]\ge0\)
for \(u=|\eta|\le1\). Hence
\(4|\eta|(1-\eta^2)/(1+\eta^2)\le2\).
Divide the terminal correction by \(E>0\), apply the sup bounds,
and add \(b_0\).
\end{proof}
The code \texttt{a11\_residual\_certificate.py} checks the
rational implication and refuses missing interval inputs.
A positive synthetic test is NOT an authenticated source certificate.
The common A.11/A.19 source roots, full-hull future-tail bound,
other exterior stress cones, Appendix B fixed point, unforced startup
systems, physical PDE residual, and the proposed unforced
Theorem 1.1 remain unproved. Lean code in this tranche covers
only finite scalar identities, not the source integrals or PDE.
